% --------------------------------------------------------------
% IEEE Article Template - LaTeX
% --------------------------------------------------------------
% Document class IEEEtran
\documentclass[conference, a4paper]{IEEEtran}

% Paquetes básicos
\usepackage[utf8]{inputenc}
\usepackage[T1]{fontenc}
\usepackage{graphicx}
\usepackage{amsmath}
\usepackage{amssymb}
\usepackage{booktabs}
\usepackage{cite}
\usepackage{xcolor}

% --------------------------------------------------------------
% Title Block - IEEE Style
% --------------------------------------------------------------
\title{\LARGE\bfseries Título del Artículo en Formato IEEE}

% Autor con afiliación
\author{
   ~\IEEEauthorblockN{Autor 1\IEEEauthorrefmark{1}}
    \IEEEauthorblockA{\IEEEauthorrefmark{1}Universidad, Departamento\\
    ciudad@universidad.edu}
    \and
    ~\IEEEauthorblockN{Autor 2\IEEEauthorrefmark{2}}
    \IEEEauthorblockA{\IEEEauthorrefmark{2}Universidad, Departamento\\
    autor2@universidad.edu}
}

\begin{document}

\maketitle

% --------------------------------------------------------------
% Abstract - 150 a 250 palabras
% --------------------------------------------------------------
\begin{abstract}
Este artículo presenta una investigación sobre [tema]. El estudio aborda [problema]
mediante [método]. Los resultados demuestran que [hallazgo principal]. La metodología
empleada consistió en [descripción]. Las conclusiones sugieren que [implicación].
\end{abstract}

% Keywords - IEEE style
\begin{IEEEkeywords}
Palabra clave 1, Palabra clave 2, Palabra clave 3, keyword4, keyword5
\end{IEEEkeywords}

% --------------------------------------------------------------
% Introducción
% --------------------------------------------------------------
\section{Introducción}
\label{sec:introduccion}
La introducción debe presentar el problema de investigación, el contexto
teórico y los objetivos del estudio~\cite{referencia1}.

\subsection{Contexto}
El contexto del problema se establece en esta sección.

\subsection{Objetivos}
Los objetivos del estudio son:
\begin{itemize}
    \item Objetivo general
    \item Objetivo específico 1
    \item Objetivo específico 2
\end{itemize}

% --------------------------------------------------------------
% Marco Teórico
% --------------------------------------------------------------
\section{Marco Teórico}
\label{sec:marco}
La revisión de literatura debe ser actual y relevante~\cite{referencia2}.

\subsection{Antecedentes}
Se presentan los antecedentes del estudio.

\subsection{Bases Teóricas}
Las bases teóricas que sustentan la investigación.

% --------------------------------------------------------------
% Metodología
% --------------------------------------------------------------
\section{Metodología}
\label{sec:metodologia}
La sección de metodología describe el enfoque de investigación.

\subsection{Diseño de Investigación}
Se describe el diseño implementado.

\subsection{Población y Muestra}
La población y la técnica de muestreo utilizada.

\subsection{Técnicas de Recolección de Datos}
Las instrumentos de recolección de datos.

% --------------------------------------------------------------
% Resultados
% --------------------------------------------------------------
\section{Resultados}
\label{sec:resultados}
Los resultados se presentan de forma clara y objetiva.

\subsection{Análisis Descriptivo}
\begin{table}[htbp]
\centering
\caption{Resultados del Análisis Descriptivo}
\label{tab:resultados}
\begin{tabular}{lcc}
\toprule
Variable & Media & DE \\
\midrule
Variable 1 & 0.00 & 0.00 \\
Variable 2 & 0.00 & 0.00 \\
\bottomrule
\end{tabular}
\end{table}

\subsection{Análisis Inferencial}
Los resultados del análisis inferencial.

% --------------------------------------------------------------
% Discusión
% --------------------------------------------------------------
\section{Discusión}
\label{sec:discusion}
La discusión interpreta los resultados a la luz de la teoría~\cite{referencia3}.

% --------------------------------------------------------------
% Conclusión
% --------------------------------------------------------------
\section{Conclusión}
\label{sec:conclusion}
Las conclusiones del estudio son:

\begin{itemize}
    \item Conclusión 1
    \item Conclusión 2
    \item Conclusión 3
\end{itemize}

% --------------------------------------------------------------
% Referencias - Formato IEEE
% --------------------------------------------------------------
\section{Referencias}
Las referencias se numeran correlativamente en orden de aparición.

\begin{thebibliography}{00}

% Artículo de revista
bibitem{referencia1}
A.~Autor, ``Título del artículo,'' \emph{Revista}, vol. 1, núm. 1, pp. 1--10, 2024.

% Libro
bibitem{referencia2}
A.~Autor, \emph{Título del libro}, 2.$^a$ ed. Ciudad: Editorial, 2023.

% Ponencia en conferencia
bibitem{referencia3}
A.~Autor, ``Título de la ponenci,'' en \emph{Proc. Conferencia}, 2023, pp. 1--5.

% Tesis
bibitem{referencia4}
A.~Autor, ``Título de la tesis,'' Ph.D. dissertation, Univ., Ciudad, 2023.

% Página web
bibitem{referencia5}
A.~Autor. (2024). Título de la página. [En línea]. Disponible: https://ejemplo.com

\end{thebibliography}

% --------------------------------------------------------------
% Anexo (opcional)
% --------------------------------------------------------------
\section{Anexo}
\label{anexo}
Información adicional.

\end{document}